\phantomsection
\section*{Ethics and Availability}\label{sec:ethics}
\addcontentsline{toc}{section}{Ethics and Availability}

All attacks were run locally against the documented, intended use of open-source code, with placeholder data and
mock resources. No real service, account, credential or person was contacted, and no third-party system was
tested in production. The study involved no human subjects: every approver is a simulated oracle or a language
model, which is why the paper makes no claim about human behavior. Language models were queried through a
commercial routing service on the authors' account, with a spending cap enforced by the experiment runner, and the
prompts contain no personal data.

The findings that concern third-party projects are the AP2 sample card and budget accumulation, rendering and
binding in MCP form-mode elicitation, and post-approval substitution in the latest Agno release. We did not
contact the maintainers. The substitution class in Agno was reported publicly by \citet{arunkumar2026loopjacking};
the AP2 observations concern sample code and integrator documentation; and the MCP behavior follows the
specification, which already directs sensitive interactions to its URL mode. The paper describes all of them as
design observations and gives no exploit code for deployed services.

The gateway, web authenticator, Tamarin models, audit and attack probes, experiment harness, raw results and the
script that generates every figure, table and in-text number (\cref{app:repro}) are available at
\codeurl{}. The scripted and AgentDojo experiments were run from a clean tree at commit \code{a0f304a}. The
language-model approvers were run at commit \code{f358b1a}, and the attack matrix at commit \code{1138121}; both
trees had uncommitted changes outside the gateway code, which the result files record.
